\section{Recomendaciones}

A partir de los resultados obtenidos y de los aspectos que permanecen pendientes, se plantean las siguientes recomendaciones para continuar fortaleciendo el desarrollo de TravelDLS.

\begin{itemize}
	\item Completar la implementación de Google OR-Tools y compararlo con Vecino Más Cercano y 2-opt utilizando las mismas condiciones, considerando principalmente la distancia total y el tiempo de cálculo.
	
	\item Fortalecer el seguimiento mediante la integración de Mapbox y el almacenamiento de ubicaciones sucesivas, de manera que el sistema pueda mostrar no solo la última posición disponible, sino también el historial del recorrido.
	
	\item Integrar una pasarela de pago en un entorno de pruebas y comprobar correctamente pagos aceptados, rechazados o duplicados, además del cálculo del porcentaje destinado al mantenimiento de TravelDLS.
	
	\item Ampliar las pruebas hacia flujos completos que involucren frontend, backend, base de datos y servicios externos, registrando los resultados e incidencias encontradas durante cada caso.
	
	\item Validar la plataforma con empresas transportistas, conductores y clientes, con el propósito de comprobar si las funcionalidades desarrolladas responden adecuadamente a las condiciones reales de uso.
	
	\item Mantener actualizados los diagramas, la documentación técnica y los manuales para conservar la correspondencia entre el diseño del sistema, las reglas de negocio y las funcionalidades implementadas.
\end{itemize}